\chapter[Sermon 55. He Exhibits a Noble Instrument of the Dragon to Be Seen, the Old Roman Empire, Which Describes What Manner a One It Is, Etc.]{He Exhibits a Noble Instrument of the Dragon to Be Seen, the Old Roman Empire, Which Describes What Manner a One It Is, Etc.}
\label{sermon:055}
\markright{Sermon 55. He Exhibits a Noble Instrument of the Dragon to Be ...}

And I saw a beast rise out of the sea, having seven heads and ten horns: and upon his horns ten crowns, and upon his head the names of blasphemy. And the beast which I saw was like a leopard, and his feet were as the feet of a bear, and his mouth as a lion. And the dragon gave him his power, and his throne, and great authority: and I saw one of his heads as it were wounded to death, and his deadly wound was healed, and all the world wondered at the beast, and they worshipped the dragon which gave power to the beast.

John proceeds to describe, through the revelation of Jesus Christ, the prominent agents of the devil, by whom he has afflicted the church of Christ with continual and most grievous persecution. And he speaks of the old and the new Roman Empire. John could not, without exceeding great danger, utter, much less describe, these things without some human assistance, and while banished and driven into exile. For the Roman Empire was considered godly, invincible, most sacred, and everlasting. Nevertheless, the Apostle speaks and writes of this in such a way that it seems he cannot avoid being considered a seditious person, and offending against the holy majesty both of the emperor and the Empire. But what, I pray, would you do, if God commanded you to speak and write so?

The world also rages at this day, when they hear kingdoms and governments rebuked by God’s word for sin and wickedness committed. Some princes, acting imperiously, issue proclamations commanding that such things be heard no more. But the Lord says in the gospel: “If these remain silent, the stones will cry out,” signifying completely that the truth must be preached, and that it cannot be suppressed or extinguished by any decrees, threats, force of arms, or punishments. Therefore, if those to whom the office of preaching is committed were to remain silent today, the Lord will stir up other preachers, who, even if the entire world says no, will bear witness to the truth. Therefore, I would counsel princes not to waste their efforts in vain attempts against God’s truth, for they shall not prevail. The truth will triumph. For he who then equipped John against the Roman Empire, at a time when it was most flourishing and powerful, the same also at this day reveals his truth to the world, now broken and aged, and will surely overcome. Woe to those stubborn natures that love to deceive. Let all preachers learn, by the example of the Apostle John, to freely utter the things they have received in commandment, and to fear no man. He who is in us (as the same John said in 1 John 4) is greater than he who is in the world.

\index{Jerome, Saint}\index{Justin Martyr}\index{Daniel (Book of)}\index{Rome}And the beast he calls the Roman Empire of great authority, and as it were godly, not without weighty considerations. For the Lord still keeps the language of Scripture, imitating Daniel, who in chapter 7 attributes the name of beast to the Roman Empire. And Saint Jerome, explaining the prophecy of Daniel, understands the beast to be the Roman Empire, and supposes that is why it is not called a lion, or a bear, or a leopard, but a beast: so that whatever cruelty anyone can imagine in beasts, by that they may understand the Romans; for surely in their manners they have shown themselves beasts. Mithridates, the renowned king of Pontus, speaking of the Romans in the thirty-eighth book of Justin, as they themselves report, says that as their founders were nurtured by sucking a wolf, so all that people have wolfish minds, never satisfied with blood, hungry and empty for rule and riches. And now how filthy beasts many Roman princes have been, their own writers testify, chiefly Suetonius, and others who have written of the lives of the emperors. And that the people of Rome were also of beastly manners, chapter 1 of the Epistle to the Romans proves.

\index{Antichrist}You will say, I know well, sins that John includes under this image refer to the entire Roman Empire. Shall we call Constantine, Theodosius, and other godly emperors beasts? I say that the Scriptures use this manner of speaking, and by beasts, in truth, understand empires, even though they do not call all those who dwell in those empires beasts without distinction. Therefore, we understand that those living a life acceptable to God are exempted within all empires, and know assuredly that neither Daniel nor John would have defiled such innocent and praiseworthy men with words. Indeed, in all this treatise concerning the empire and Antichrist, we always except those who are innocent and excel in virtue. Of this we shall happily speak more hereafter.

And first he shows the beginning of this Empire. A beast comes out of the Sea, on the sand whereof stands the Dragon: and in the 17th Chapter it is said how the beast came out of the bottomless pit. Therefore the beginning of it is referred to Satan. Nevertheless, we must here take diligent heed that we do not take anything away from the Lord our God, which he claims for himself. Scripture in sundry places, but chiefly by two most excellent witnesses, by Daniel in the 3rd Chapter and Paul in the 13th Chapter to the Romans, has set forth that kingdoms and empires are of the Lord, and that he sets up and deposes kings. There is no power, says the Apostle, but of God. And hitherto the apostles command obedience to princes and magistrates. How then is it that we hear that the Roman Empire came out of the bottomless pit, since the Apostle speaks of the same? Doubtless the Roman Empire is not absolutely of the Devil. For God is the author of monarchies, and preserves realms and policies, giving thereunto certain faithful servants. But Satan meddles with men's matters and corrupts both kings and kingdoms; and so long they are of the Devil. The Christians in all political matters obeyed emperors, but commanding idolatry they obeyed them not. Certain it is that God did institute the kingdom of Israel or of ten tribes by the prophet Ahab; yet nevertheless the Lord cries out in another prophet, they have reigned indeed, but not by me. For the Lord would have had those kings to have framed all things after his word, and to reign in the fear of God: and where they did not so, but following the instigation of Satan ordered all things after their own lust, they are rightly said to reign, not of God, but of the devil. Therefore have the godly obeyed kings: but they obeyed them not commanding wicked things, although they took them for their kings. God had instituted the order of priests; nevertheless Christ calls the doings of the same priests the works of darkness. And Peter says: we must rather obey God than men. So verily the Roman Empire, which was of God, came also out of the Sea (as Daniel says also) out of the troublesome world, and even out of Hell, being made great through slaughter, murder, sedition and treason. For the people of Rome with the most part of emperors regarded the Devil and the world, and not God.

And what the Roman empire is today, he now depicts: it has seven heads and ten horns, and each horn has a crown, signifying truly that horns represent kingdoms. We are not here introducing any new or strained interpretation. In chapter 17, the angel himself explains this, saying that the seven heads signify seven mountains or hills, and also kings. Rome is considered to have many hills, but there are seven notable ones: the Palatine, Capitoline, Aventine, Caelian, Esquiline, Viminal, and Quirinal. Propertius explains the same in one verse (which I have expressed in two): “Seven hills the city rises above, overlooking the entire world.”

\index{Apocalypse (Book of)}And therefore is called by the Greeks, “Apocalypse,” meaning “revelation” or “discovery.” And indeed, the city is taken to represent the entire Empire. So there have been many kings and emperors encompassed within the seventh number. But it is certain that the seventh number of kings is exactly found in the history. For at the beginning, when Rome was first built, there reigned seven kings in succession: Romulus, Numa, Tullus Hostilius, Ancus Marcius, Tarquinius Priscus, Servius Tullius, Tarquius Superbus. They were expelled because Lucretia was ravished by the king’s son, and they were ruled by consuls, by ten men, and by dictators, until the time of Julius Caesar, who first usurped a kingly crown for himself. After him reigned Antony and Octavian, called Augustus, Tiberius, Caligula, Claudius, and Nero, again seven. In Nero, the Empire receives a plague. From thence again are accounted seven: Otho, Galba, Vitellius, Vespasian, Titus, Domitian, Nerva. From him, the Empire devolved to Ulpius Trajan, a Spaniard. Therefore, the Roman Empire could not be expressed by plainer marks. To this Empire also Daniel attributed ten horns, both because it was composed of many kingdoms and also because it was dispersed again into many. This will be spoken of in chapter 17. And it is a common thing in the Scriptures to signify kingdoms and power by horns.

And to this kingdom the Lord Jesus ascribes open wickedness, which he calls blasphemy. For he adds: and upon his heads the name of blasphemy, that is to say, whatever blasphemy may at any time be devised anywhere will be found manifest in this Empire, and chiefly in its leaders. For if one beholds the hills of Rome, chiefly the Capitoline Hill, one will find it called by Cicero the dwelling place of the gods, truly because it contained, in a manner, the images of all the gods. For on those hills were seen the temples of Jupiter with all his attributes, and so forth; the temples of Saturn, Juno, Minerva, Mars the avenger, Hercules, Janus, Venus, Apollo; also the temples of Fortune, Health, Victory, Concord, and such others. But if one looks upon the emperors themselves, Gaius wanted his images set up in temples, and the people to swear by his name. Nero blasphemed the name of Christ, and by shedding innocent blood sought to abolish the Gospel. Domitian commanded himself to be called God and Lord. And others also have demanded divine honors, men steeped in blasphemies, and reeking in all wickedness.

Furthermore, by an image composed of diverse beasts, he shows how the Roman Empire increased and obtained such power, and what its manners were. In the seventh chapter of Daniel, the eating of the mountain signifies the monarchy of Greece or Macedon, the bear signifies the Persian monarchy, and the lion signifies the monarchy of the Chaldeans or Babylonians. And it is plain that the Romans, overcoming those nations and subjugating those monarchies to themselves, came to the supreme height of governance. For they subdued the eastern regions chiefly by Lucullus, Pompey, and Crassus; Macedon and all Greece by Paulus Aemilius; a good part of Africa by Scipio and Marius; Egypt by Octavius Augustus; and so forth. And just as they were in religion ungodly, so in other manners they were not unlike wild beasts. For as the Libyan or panther is spotted and of diverse colors, so are the Romans, a collection of many nations, prone to sedition and slaughter. The bear only goes upon his feet, but with them also strikes and catches his prey; so the Romans did nothing else but strike, fight, and take spoils. And just as the strength of a lion is most excellent among four-footed beasts, and the lion’s mouth is insatiable and stinking, so was the Roman Empire most strong, covetous, never contented, and the very essence and corruption of mischief.

\index{John, Saint}And Saint John declares indeed more expressly, that the Romans possess from the Devil all wickedness, cruelty, and mischief. The Dragon says, “He gave to that beast his power, and that great authority.” This is, in effect, as if he had said: The Devil reigned wholly in the Romans, and the Romans wrought by the Devil all that they did. For the Devil is the origin of murders and lies. Of the Devil’s seat I have spoken in the second chapter of this book. However, we must know that all power is of God; but he, by his just judgment, permits many things to the Devil over the children of unbelief. For Saint Paul in 2 Thessalonians 2, when he had spoken of the mighty working of Satan, by signs and false wonders, with which they should be deceived, those who would not receive the truth, he adds immediately: “Therefore God shall send them strong delusions, that they may believe lies, and be judged all those who believed not the truth,” and so forth. For (as I have often admonished) we must take good heed that we mix not the works of God and the Devil together. Good works are of God, evil are of the Devil. Now, lest any man should marvel why God permits so much to the Romans and the Devil their head, and does not infringe their force for the elect’s sake, Saint John interweaves a heavy fate of the people of Rome and of the whole Empire, which befell them immediately after the first persecution raised against the church of Christ, and after the most shining Apostles were executed, verily to avenge that innocent blood. For he sees one of those heads, as it were, wounded to death. For Nero, who first of the Emperors, stirred up the first persecution against the church, killed himself. And he was the last Emperor of that family. And he left the Empire so afflicted that it was likely to fall to decay. Certain provinces revolted. Galba, Otho, and Vitellius fought among themselves and made civil wars. This Vitellius moreover, drove Sabinus, Vespasian’s brother, suspecting no evil, with others, into the Capital house, and setting the temple on fire, destroyed both Temple and men together, and made all one heap. Neither does Orosius conceal why these things happened, saying: “Rome fell swiftly by the murder of princes and civil wars, for the injuries done to the Christian religion.”

Nevertheless, the Apostle adds that the wound was healed again. For Sextus Aurelius Victor says that Vespasian, in a short time, refreshed the weary world that had been wounded. Here, he says, is the deadly wound that was healed again. For other writers, discussing this matter more fully, describe how Vespasian, returning to Rome, considered nothing more noble or better than to establish and beautify the commonwealth that was sorely afflicted and decayed, to bring order to the provinces and cities that were disordered by tumults and seditious uprisings, to reform the military discipline that had become lax, and to punish offenders. He repaired the city, damaged by old fires and ruins, with new buildings; he rebuilt the Capitol that had been burned; and he erected the Theater in the midst of the city, the most ancient monument of the Empire, and so on.

Moreover, he now touches sharply upon the foolishness and wickedness of the world. And there was great wonder throughout the earth, and so forth. For the world follows present fortune, and judges all things by their good or evil outcomes. For that religion, they say, is most noble, stable, and true, which is celebrated by victories and shines with the ornaments of this world. Therefore, because of the majesty of the Roman Empire, which they held in great admiration, most people received the Roman religion and defended it as sincere. But John, declaring the enormity of this sin, says, “And they worshipped the dragon,” and so forth. He does not say that they worshipped gods, or wood and stones, but that they worshipped the Devil. Idolaters will say that they worship and honor gods, and are not ignorant that images are made of corruptible material? And that the worship they give to them returns, not to those dead signs, but to those of whom they are signs. Thus, truly, all idolaters say. And if you say to them, “You worship wood and stones,” they will quickly answer that great injury is done to them. For they are not so foolish, they will say, as to worship that which they made with their own hands, and so forth. But the Apostle, who knew well these civil explanations and crafty shifts of idolaters, speaking frankly against them, did not respect what they alleged for themselves, but rather what God judges, and the truth of the thing declares, and says, “And they worshipped the Devil or the Dragon.” So Paul in 1 Corinthians, chapter 10. “The things that the Gentiles offer up,” he says, “they offer to demons.” But the Gentiles deny this. But God in this case does not pass judgment on the judgments, intentions, and denials of men, but pronounces according to his own judgment. In Leviticus, chapter 17, he says, “If anyone offers oblations to me other than those I have prescribed, they shall defile yourselves with blood.” Let now the Mass-singing priests cry out until they are hoarse again, “We offer to the Lord God, not to strange gods!” Yet the Lord’s sentence shall stand most true forever, that they transgress with unlawful worship no less than if they committed parricide. As also Isaiah bears witness in chapter 66. The Lord God allows the sincere obedience that we show to his laws; he cares nothing for our inventions and good intentions. Thus, at this present, he shows in few words what the thing is in truth, that all idolaters worship the Devil. If we would at this day esteem these things rightly, we should not contend as if for life and lands, about maintaining images in the church. The Lord Jesus enlighten our hearts and minds to see his truth.
